\documentclass{report}
\usepackage{amsmath,amssymb}
\setlength{\parindent}{0mm}
\setlength{\parskip}{1em}
\begin{document}
\begin{center}
\rule{6in}{1pt} \
{\large Keiichi Morikuni \\
{\bf Statioary inner-iteration preconditioned GMRES method for least squares problems}}

The Graduate University for Advanced Studies \\ 2-1-2 Hitotsubashi \\ Chiyoda-ku \\ Tokyo 101-8430 Japan
\\
{\tt morikuni@nii.ac.jp}\\
Ken Hayami\end{center}

\newtheorem{theorem}{Theorem}
\newtheorem{definition}{Definition}

Consider solving large sparse linear least squares problems
\begin{align}
\min_{\boldsymbol{x} \in \mathbf{R}^n}{\left\| \boldsymbol{b} - A
\boldsymbol{x} \right\|_2},
\label{eq:1.1}
\end{align}
where $A \in \mathbf{R}^{m \times n}$, $\boldsymbol{b} \in \mathbf{R}^m$,
and $A$ is not necessarily of full rank.
The least squares problem \eqref{eq:1.1} is equivalent to the following normal equation
\begin{align}
A^\mathsf{T} A \boldsymbol{x} = A^\mathsf{T} \boldsymbol{b}.
\label{eq:1.2}
\end{align}
The standard iterative methods are the (preconditioned) CGLS [M. R.
\break Hestenes and E. Stiefel, J. Research Nat. Bur. Standards, 49
(1952), pp.~409--436] and LSQR [C. C. Paige and M. A. Saunders, ACM
Trans. Math. Software, 8 (1982), pp.~43--71] methods.
However, since $\kappa_2 (A^\mathsf{T}A) = {\kappa_2 (A)}^2$, iterative
methods may be slow to converge.
Here, $\kappa_2 (A) = \sigma_\mathrm{max} / \sigma_\mathrm{min}$ is the condition number,
where $\sigma_\mathrm{max}$ and $\sigma_\mathrm{min}$ are the largest and
smallest positive singular values of $A$.
Hence, when iterative methods are used, good preconditioners are
necessary to achieve better convergence.

In this paper, we propose using stationary inner-iteration
preconditioners for the generalized minimal residual method (GMRES) for
solving least squares problems, and present theoretical justifications
for using the inner iteration as preconditioners even for rank-deficient
matrices.
For the preconditioners, the Jacobi overrelaxation and successive
overrelaxation (SOR)-type iterative methods designed for least squares
problems are employed.
These inner iterations are efficient in terms of computational work and memory.
A significant advantage of such inner-iteration preconditioners is that
one can avoid explicitly computing and storing the preconditioning
matrix.
Numerical experiments for large sparse least squares problems including
ill-conditioned and rank-deficient ones show that the proposed methods
outperform previous methods.

The main motivation for proposing the new preconditioners is to reduce
CPU time and memory significantly, and to broaden the scope of problems
that can be solved.
Note that, previously, GMRES preconditioned by the robust incomplete
factorization (RIF) [M. Benzi and M. T{\r u}ma, SIAM J. Sci. Comp., 25
(2003), pp.~499--512] was comparable with, but not definitely superior
to, reorthogonalized CGLS with RIF in terms of CPU time for
ill-conditioned problems [K. Hayami, J.-F. Yin and T. Ito, SIAM J. Matrix
Anal. Appl., 31 (2010), pp.~2400--2430].
Moreover, RIF for least squares problems will break down for rank-deficient matrices.
We aim to improve on these points.

The left-pre\-con\-di\-tioned GMRES method for least squares problems
proposed in [Hayami \textit{et al.}, 2010] is called the BA-GMRES method.
It applies GMRES to $\displaystyle \min_{\boldsymbol{x} \in
\mathbf{R}^n}\| B \boldsymbol{b} - B A \boldsymbol{x} \|_2$, where $B \in
\mathbf{R}^{n \times m}$ denotes the preconditioning matrix.
Conditions for the preconditioner $B$ for BA-GMRES were given in [Hayami
\textit{et al.}, 2010].

Instead of explicitly forming and storing an preconditioning matrix $B$,
we propose applying several steps of a certain stationary iterative
method for solving the normal equation \eqref{eq:1.2} of the form
\begin{align}
\boldsymbol{x}^{(l+1)} = M^{-1} N \boldsymbol{x}^{(l)} + M^{-1}
A^\mathsf{T} \boldsymbol{b}, \quad l = 0, 1, \dots, \quad
\boldsymbol{x}^{(0)} = \boldsymbol{0},
\label{eq:2.1}
\end{align}
inside BA-GMRES, where $A^\mathsf{T} A = M - N$, $M$ is nonsingular, and
$l$ is the number of inner iterations.
Define the iteration matrix by $H = M^{-1} N$.
Tanabe [Numer. Math., 22 (1974), pp.~349--359] defined the following.
\begin{definition}
$H$ is convergent if $\displaystyle \lim_{l \rightarrow \infty} H^l$
exists, and divergent otherwise.
\end{definition}

Then, we obtain the following theorem for the stationary inner-iteration BA-GMRES method.
\begin{theorem} \label{th:2.6}
Assume that $H$ is convergent.
Then, BA-GMRES with the inner-iteration preconditioning of the form
\eqref{eq:2.1} determines a least squares solution of $\displaystyle
\min_{\boldsymbol{x} \in \mathbf{R}^{n}} \| \boldsymbol{b} - A
\boldsymbol{x} \|_2$ without breakdown for all $\boldsymbol{b} \in
\mathbf{R}^m$.
\end{theorem}
According to [A. Dax, SIAM Review, 32 (1990), pp.~611--635], the
Cimmino-NR and NR-SOR methods [Y. Saad, 2nd ed., SIAM, Philadelphia,
2003] give a convergent iteration matrix so that these methods can be
used as the inner iterations for BA-GMRES.

Moreover, we show that the stationary inner-iteration BA-GMRES residual
bound depends on the spectral radius of $H$.

Finally, we will present numerical results that show BA-GMRES
preconditioned by the NR-SOR inner iteration outperforms conventional
methods for ill-conditioned, rank-deficient, and practical problems.
A strategy for tuning the parameters for the NR-SOR inner iteration will
be also proposed and shown to be effective.


\end{document}
